\section{SF-E827. Diagonal Yoneda classes survive arbitrary cross-index
changes}\label{SF-E827:sf-e827.-diagonal-yoneda-classes-survive-arbitrary-cross-index-changes}

Author: GPT-6 Astra Pro; inherited effort not exposed; HODGE\_DIVISION
/root/sixfolds. Date: 2026-11-21.

\textbf{Status:} independently whole-reviewed by /root and
/root/g34\_eightfolds; both mathematical reviews PASS on the final frozen
draft
\texttt{abe42293fb48da62436bfe7fafd7ae19a2546a768ad294f6f89ed4ec451f8a27}.
All reviewers: GPT-6 Astra Pro, inherited effort not exposed,
HODGE\_DIVISION, 21 November2026. No new numerical computation is a premise.
This is a separate strengthening of E826. It removes the global Ext3(B,A)=0
hypothesis in specified line-class settings by following diagonal classes
through the actual triangle maps. It does not claim all positive-determinant
objects or all divisor directions are excluded.

\subsection{1. Exact setting and a diagonal-survival
lemma}\label{SF-E827:exact-setting-and-a-diagonal-survival-lemma}

Let X be an abelian fourfold. Consider an actual coherent extension

\[
 0\to A\xrightarrow{i}E\xrightarrow{p}B\to0,\qquad e:B\to A[1],
 \quad
 A=A_0\oplus\bigoplus_{a=1}^{u}L_a^{\oplus n_a},
 \quad B=\bigoplus_{b=1}^{v}M_b^{\oplus m_b}.
 \tag{E827.1}
\]

Assume a specified coherent finite effective translation/tensor action
preserves this exact sequence and the displayed blocks. These are actual
equivariant summands of A and B, not asserted summands of E. A0 is allowed
to be zero or an arbitrary coherent perfect object in this lemma. In the
application A0 is the unit block and the intrinsic torsion/HN argument of
E826 supplies all the assumed equivariance for every action on E.

Put V=H1(X,O\_X), so H2(O\_X)=Lambda2 V. Let W,Z be vector subspaces of
H2(O\_X) with W intersect Z zero. Assume for every a,b:

\begin{enumerate}
\def\labelenumi{\arabic{enumi}.}
\tightlist
\item
  the Yoneda product of Ext1(L\_a,M\_b) with Ext1(M\_b,L\_a) has image in Z
  inside Ext2(L\_a,L\_a)=H2(O\_X), and its reversed product has image in Z
  inside Ext2(M\_b,M\_b);
\item
  multiplication by W kills Ext1(M\_b,L\_a) into Ext3(M\_b,L\_a), and also
  kills Ext1(L\_a,M\_b) into Ext3(L\_a,M\_b).
\end{enumerate}

Let J be any subset of the B indices such that Ext3(M\_b,A0)=0 for b in J.
Then

\[
 \operatorname{ext}^2_\Gamma(E,E)\ge(\dim W)(u+|J|).
 \tag{E827.2}
\]

If A0=0, J may consist of all v indices. The statement also holds without an
action, by taking the trivial group. No vanishing of all cross Ext3 groups
is assumed.

Here is the complete proof, retaining multiplicity matrices. In Ext2(A,A)
let D\_A be the direct sum over a of W tensor the identity matrix on the
multiplicity space of L\_a. Thus dim D\_A=(dim W)u. Its elements have zero
A0 block and zero off-diagonal blocks. Each is invariant: translations and
degree-zero tensoring act trivially on H2(O\_X), and conjugation fixes the
identity of any multiplicity representation, including a projective one.

In Hom(A,triangle), the kernel of
\(i_*:\operatorname{Ext}^2(A,A)\to\operatorname{Ext}^2(A,E)\) is the image
of Ext1(A,B). For a class gamma in that group, its image is \(e[1]\gamma\).
Its a-th diagonal line block is exactly

\[
 \sum_b e_{ab}[1]\gamma_{ba},
 \tag{E827.3}
\]

with the multiplicity matrices multiplied in this order. Hypothesis1 places
every coefficient in Z, hence the whole block in Z tensor
End(C\^{}\{n\_a\}). This has zero intersection with W tensor its identity.
Thus D\_A injects into Ext2(A,E), regardless of any off-diagonal image or
any component involving A0. The absence of A0 in B is used here.

In Hom(-,E), the next boundary from Ext2(A,E) to Ext3(B,E), on a class
\(i[2]d\) with d in D\_A, is \(i[3]d[1]e\). Its only nonzero possible
components before postcomposition by i are \(d_a[1]e_{ab}\). They are zero
by hypothesis2, also with arbitrary multiplicity matrices, because d\_a is a
scalar identity times a class in W. Consequently the image I of Ext2(E,E) in
Ext2(A,E) contains this injected D\_A.

The B-side contributes independently to the kernel of that map. In
Ext2(B,B), let D\_B be the sum over b in J of W tensor the corresponding
identity matrices. The connecting map from Ext2(B,B) to Ext3(B,A) sends such
a class h to \(e[2]h\). Every line-target component is zero by hypothesis2,
and every A0 component is zero by Ext3(M\_b,A0)=0. Hence D\_B lifts to
Ext2(B,E).

We must check that the preceding boundary in Hom(-,E) cannot kill these
lifted diagonal classes. For any alpha in Ext1(A,E), its boundary is
\(\alpha[1]e\). Projecting to Ext2(B,B) gives

\[
 (p[1]\alpha)[1]e.
 \tag{E827.4}
\]

The b-th diagonal block, b in J, is a sum of line-to-line degree-one
products in Z tensor End(C\^{}\{m\_b\}). The additional possible term
through A0 is zero: by fourfold Serre duality Ext1(A0,M\_b) is dual to
Ext3(M\_b,A0), which is zero. Therefore the diagonal projection of every
such boundary has no W component after passing to H2(O\_X)/Z.

More formally, take diagonal projection, matrix trace divided by m\_b, and
the quotient by Z on each selected block. This is an equivariant linear map
from Ext2(B,B) to a space containing an injected copy of
\(\bigoplus_{b\in J}W\). It kills the projection of all the boundaries
(E827.4), while the lifts of D\_B map onto that copy. Thus

\[
 Q=\operatorname{Ext}^2(B,E)\big/\operatorname{im}
     \bigl(\operatorname{Ext}^1(A,E)\to\operatorname{Ext}^2(B,E)\bigr)
 \tag{E827.5}
\]

has dimension at least(dim W)\textbar J\textbar, and injects into Ext2(E,E).
Its quotient there is I, which already contains D\_A. All constructions
commute with the finite action. Taking invariants is exact; invariant lifts
of D\_B exist by averaging; W and the identity matrices remain fixed. The
same dimension bounds therefore hold for Q invariants and I invariants. This
proves (E827.2). These are the two actual triangle boundary maps, not an
assumed associated-graded decomposition of E.

\subsection{2. Nondegenerate line differences: six classes per positive
block}\label{SF-E827:nondegenerate-line-differences-six-classes-per-positive-block}

Suppose every line bundle L\_a tensor M\_b inverse is nondegenerate. The
abelian index theorem puts all its cohomology in a single degree j, and that
of its inverse in degree4-j. Thus its Ext1 and the reverse Ext1 cannot both
be nonzero. Their products in either diagonal self-Ext2 are zero. Moreover,
whenever its Ext1 is nonzero, j=1 and its Ext3 is zero; the same applies to
the inverse. Hypotheses1--2 of section1 hold with

\[
 Z=0,\qquad W=H^2(X,O_X),\qquad\dim W=6.
 \tag{E827.6}
\]

Consequently ext2\_Gamma(E,E) is at least6(u+\textbar J\textbar). In
particular u at least4 already excludes a weakly semiregular balanced
character of cap rank18, without \textbf{any} condition on Ext3(B,A0) or on
other cross Ext3 groups. Opposite index-three pairs elsewhere do not cancel
the diagonal classes: the extension component and the relevant reverse
degree-one component would have to occur in the same nondegenerate pair,
where they cannot.

For the two-slope architecture in E826, its canonical torsion/HN filtration
and pairwise block Hom vanishings provide the equivariance required in
section1 for every specified finite action on E. Thus this is a theorem for
all extensions in that architecture satisfying the nondegenerate
line-difference condition, not merely for the printed E826 example.

\subsection{3. Degenerate differences in the four-divisor Gaussian
algebra}\label{SF-E827:degenerate-differences-in-the-four-divisor-gaussian-algebra}

There is a stronger conclusion for the actual divisor algebra used in our
constructions, even when some cross differences are degenerate. Write

\[
 X=Y\times Z_0=E_i^2\times E_i^2,
 \qquad V_Y=H^1(O_Y),\quad V_Z=H^1(O_{Z_0}).
 \tag{E827.7}
\]

These are the two S-matrix factors; they are \textbf{not} the irrational RM
embedding planes. In the divisor basis(theta,theta\_U,theta\_S,theta\_US), a
rational difference has two blocks xI+yU2, with U2 squared equal to2. A
nonzero such block is nondegenerate: x2-2y2=0 with rational x,y forces
x=y=0. Hence a nonzero difference is either nondegenerate on X, or is a
pullback of a nondegenerate line class from exactly one of Y,Z0. The
opposite theta slopes in E826 ensure that a cross difference is nonzero.
Arbitrary degree-zero twists do not change the following proof.

Set

\[
 W=V_Y\otimes V_Z,\qquad
 Z=\Lambda^2V_Y\oplus\Lambda^2V_Z
       \ \subset\ H^2(O_X),\qquad\dim W=4.
 \tag{E827.8}
\]

This is a direct decomposition H2(O\_X)=W plus Z. The two projections are
fixed by translations and degree-zero tensor operations on cohomology. We
verify all hypotheses of section1, including every degenerate case. The
nondegenerate case has already been proved in section2.

Suppose the difference line is pulled back from Y. It is an exterior product
L\_Y box P\_Z, up to a degree-zero twist on Y, because Pic0 of a product is
the product of the dual abelian varieties. If P\_Z is nontrivial, all its
cohomology, and all that of its inverse, vanish. Otherwise P\_Z is trivial.
Let j in\{0,1,2\} be the index of the nondegenerate surface line L\_Y.
Kunneth and the surface index theorem give

\[
 H^1(X,L_Y\boxtimes O_{Z_0})=
 \begin{cases}
 H^0(Y,L_Y)\otimes H^1(O_{Z_0}),&j=0,\\
 H^1(Y,L_Y)\otimes H^0(O_{Z_0}),&j=1,\\
 0,&j=2.
 \end{cases}
 \tag{E827.9}
\]

Multiplication by W kills each of these spaces into H3 of the line: it
raises the Y cohomological degree by one, where the nondegenerate line has
no cohomology. The same statement holds for its inverse, whose index is2-j.
Both degree-one spaces can be nonzero only when j=1. In that case their
product lies in H2(O\_Y) tensor H0(O\_Z0), the first pure summand of Z. This
proves hypotheses1 and2, with all multiplicity matrices allowed. The case
pulled back from Z0 is identical with the factors interchanged. This is an
actual cohomological product calculation, not only a count of the dimensions
of the Ext groups.

Therefore for all nonzero cross differences in this rational four-divisor
algebra,

\[
 \operatorname{ext}^2_\Gamma(E,E)\ge4(u+|J|).
 \tag{E827.10}
\]

In particular u at least5 excludes every such two-slope balanced extension,
even with degenerate differences and nonzero opposite Ext3 pairs. If one
uses selected negative blocks in J, the threshold is u+\textbar J\textbar{}
at least5. The defining condition for J is still the explicit Ext3(M\_b,A0)
vanishing; no assertion about the unit block is silently made.

\subsection{4. Balanced consequence and exact remaining
scope}\label{SF-E827:balanced-consequence-and-exact-remaining-scope}

The entire normalized balanced character, with any algebraic rational B
field and q nonzero, has degree-two cap rank18. The exterior proof is
repeated in E826: two pure blocks have rank1 each and the mixed block has
rank16. The B-shear is invertible and preserves this rank. The supplied
invariant CY4 adjunction(S21.8) therefore requires ext2\_Gamma(E,E)=18 under
weak semiregularity. Thus either lower bound exceeding18 is a contradiction
for every specified coherent effective finite action in that framework.

This is a strict enlargement of E826's family exclusion: a merely nonzero
Ext3(B,A) no longer suffices to escape, and in the displayed CM divisor
algebra neither does a degenerate cross pair. The argument still requires
the actual equivariant filtration and blocks, supplied by the proved
two-slope/torsion construction in the applications. It does not assert that
an arbitrary complex or unstable sheaf has such a filtration. When there are
too few positive blocks, selected negative blocks do not meet their
additional vanishing condition, or divisor directions leave the four-divisor
algebra, this proof leaves a genuine gate. No successful seed or general
nonexistence theorem is claimed.

\subsection{5. Source and check boundary}\label{SF-E827:source-and-check-boundary}

The abelian and surface index/vanishing inputs are the original Mumford
Abelian Varieties(1970), section16 printedpp150--153 and162--163, already
read in this run. Nontrivial Pic0 cohomology vanishing is the previously
read published Mukai1981 Theorem2.2 p156 and its base-change consequence, or
the direct flat-character argument already supplied. The Kunneth
factorization, both actual triangle boundary calculations, and all
multiplicity/invariant projections are proved above. The full S21 adjunction
and its stated source statuses are retained exactly as freshly reread for
E826; no new published status is asserted for a supplied preprint input.
BF2008's exact published evaluation/trace locators remain those fully
recorded in E826/E825.

There is no new numerical-rank premise. This proof is a symbolic
Yoneda/Kunneth argument, independently whole-reviewed as recorded above. The
existing E826 exact computation supplies an actual example in its scope but
does not replace the proof for the larger family here.
